% =========================================================
% ANÁLISIS TEMPORAL 2019–2024 — versión RUCP (~800 palabras)
% =========================================================

\section{Evolución Temporal: 2019--2024}
\label{sec:temporal}

\subsection{La Paradoja de la Coalición Cohesiva Derrotada}

La comparación entre las elecciones de 2019 y 2024 revela una paradoja: la Coalición Republicana perdió el balotaje de 2024 a pesar de mejorar sustancialmente su cohesión interna respecto a 2019. La Tabla~\ref{tab:comparison_2019_2024} presenta las estimaciones mediante promedios departamentales ponderados por votos, que permiten comparabilidad entre años con cambios en la distribución de circuitos.

\begin{table}[H]
\centering
\caption{Evolución de defecciones coalicionales 2019-2024. Mejora dramática en cohesión: todos los partidos redujeron defección entre 26.6 y 92.1 puntos porcentuales.}
\label{tab:comparison_2019_2024}
\small
\begin{tabular}{lrrrrr}
\toprule
\textbf{Partido} & \textbf{Def. 2019} & \textbf{Def. 2024} & \textbf{Cambio} & \textbf{Votos def.} & \textbf{Votos def.} \\
 & \textbf{(\%)} & \textbf{(\%)} & \textbf{(pp)} & \textbf{2019} & \textbf{2024} \\
\midrule
Cabildo Abierto & 69.5\% & 27.6\% & -41.9 & 186.316 & 16.382 \\
Partido Colorado & 33.8\% & 7.2\% & -26.6 & 101.412 & 27.793 \\
Partido Nacional & 36.0\% & 3.6\% & -32.4 & 250.260 & 23.077 \\
Partido Independiente & 94.8\% & 2.8\% & -92.1 & 22.335 & 1.130 \\
\bottomrule
\multicolumn{6}{l}{\footnotesize Defección = \% de votos transferidos de partido coalicional a FA en balotaje.} \\
\multicolumn{6}{l}{\footnotesize Partido Independiente tuvo el cambio más dramático: -92.1 pp (de 94.8\% a 2.8\%).} \\
\end{tabular}
\end{table}


\textbf{Cabildo Abierto}: en 2019 el 88.7\% de los votantes de CA transfirió su voto al PN en el balotaje; en 2024, esta proporción cayó al 46.8\%. La defección al FA aumentó de 10.3\% a 52.7\%, un deterioro de 42.4 puntos porcentuales. Este colapso en la retención coalicional es el factor más determinante de la derrota de 2024.

\textbf{Partido Colorado}: la defección PC$\to$FA se \textit{redujo} de 36.3\% en 2019 a 9.7\% en 2024 ($-26.6$ pp). Esta mejora fue la más marcada de todos los partidos de la coalición, pero su impacto positivo quedó eclipsado por el deterioro de CA.

\textbf{Partido Nacional}: la defección PN$\to$FA también se redujo de 38.6\% en 2019 a 6.2\% en 2024 ($-32.4$ pp). La retención del PN mejoró de 60.2\% a 91.8\%.

\textbf{Partido Independiente}: el PI pasó de una defección al FA del 92.0\% en 2019 a virtualmente 0\% en 2024, el cambio más drástico de todos los partidos ($-92.1$ pp). Su rol cambió de mayor fuente de fugas a partido más leal de la coalición.

En suma: tres de los cuatro socios de la coalición mejoraron sustancialmente su retención. La coalición perdió igualmente. ¿Cómo es posible?

\subsection{El Factor Decisivo: Colapso del Caudal de CA}

La respuesta está en el tamaño. En 2019, CA obtuvo 268{,}134 votos en primera vuelta; en 2024, solo 59{,}412 ($-$208{,}722 votos, $-$78\%). La mejora en la tasa de retención de cada partido no alcanza a compensar que la base a retener se redujo drásticamente. La coalición ingresó al balotaje de 2024 con 154{,}000 votos menos que en 2019 para convertir a la opción de la coalición, independientemente de cuán bien performaran las tasas de retención.

La Figura~\ref{fig:temporal_sankey} ilustra el contraste: los flujos de 2019 muestran una coalición más fragmentada pero más grande; los de 2024, una coalición más cohesiva pero pequeña, con una concentración de pérdidas en el propio caudal de CA.

\begin{figure}[H]
\centering
\includegraphics[width=0.90\textwidth]{temporal_comparison_flows.png}
\caption{Comparación de flujos de transferencia 2019 vs.\ 2024. Las flechas representan proporciones estimadas de transferencias de votos entre partidos de origen y opciones de balotaje. El colapso del caudal de CA entre elecciones es la variable explicativa central del resultado de 2024.}
\label{fig:temporal_sankey}
\end{figure}

\subsection{Nula Estabilidad Geográfica}

Un resultado adicional del análisis temporal es la ausencia de correlación entre los patrones de defección departamentales de 2019 y 2024. Las correlaciones entre las tasas de defección CA$\to$FA por departamento en ambas elecciones son $r < 0.13$ (prácticamente cero). Esto implica que los departamentos con alta defección en 2019 no fueron necesariamente los mismos que en 2024, y viceversa. La reconfiguración geográfica es profunda: no existe una geografía estable de la deslealtad coalicional en Uruguay, sino respuestas locales a condiciones coyunturales específicas de cada elección.

Esta inestabilidad geográfica es relevante para la planificación electoral: no es posible predecir con base en 2019 qué departamentos serán los focos de defección en una elección futura. La heterogeneidad entre elecciones es tan grande como la heterogeneidad dentro de cada elección.
